% !TeX program = lualatex
\documentclass[a4paper,11pt]{ltjsarticle}
\usepackage{amsmath,amssymb,amsthm}
\usepackage{lmodern}
\usepackage{booktabs}
\usepackage{array}
\usepackage[unicode]{hyperref}

\theoremstyle{definition}
\newtheorem{definition}{定義}[section]
\newtheorem{assumption}[definition]{仮定}
\newtheorem{remark}[definition]{注記}
\newtheorem{example}[definition]{例}
\newtheorem{conjecture}[definition]{予想}
\theoremstyle{plain}
\newtheorem{theorem}[definition]{定理}
\newtheorem{lemma}[definition]{補題}
\newtheorem{proposition}[definition]{命題}
\newtheorem{corollary}[definition]{系}

\newcommand{\Set}{\mathbf{Set}}
\newcommand{\Hom}{\operatorname{Hom}}
\newcommand{\End}{\operatorname{End}}
\newcommand{\id}{\mathrm{id}}
\newcommand{\Hol}{\operatorname{Hol}}
\newcommand{\dfix}{\bigcirc}
\newcommand{\dline}{\mathrel{\rule[0.55ex]{1.6em}{0.4pt}}}

\title{向の理論の展望\\[0.2em]
  {\Large ---境界接続、投機、補償および残余---}\\[0.4em]
  {\large Perspectives on the Formal Theory of Direction}}
\author{千葉将臣}
\date{2026-08-26}

\begin{document}
\maketitle

\begin{abstract}
本稿は、向の形式理論を圏論以前の代替的基礎としてではなく、対象化された内部世界と、その世界では終域として所有されない境界との接続を記述する表示理論として再定位する。通常の内部遷移 $A\to B$ は、始域と終域が既に対象として与えられた後の外延的関係である。これに対して
\[
 A\dline\dfix
\]
は、現在の観点では終域を内部対象として形成できない半開方向を表す。この区別は絶対的な圏外性を主張するものではない。十分大きな拡大圏では境界方向を射として表現できるが、その場合にも元の観点への非降下性と非表現可能性が残る。

本稿の中心は、正の境界方向前層 $\partial_\omega^+$、負の境界方向余前層 $\partial_\omega^-$、および両者を内部射へ部分的に送る境界接続子
\[
 \kappa_{A,B}:
 \partial_\omega^+(A)\times\partial_\omega^-(B)
 \rightharpoonup
 \Hom_{\mathcal C}(A,B)
\]
である。対角成分は境界ホロノミーを与え、恒等帰還、非恒等な残余、帰還不能を区別する。この構造は、未確定な未来への投機、検証による帰還、補償、ロールバックおよび副作用の痕跡を統一的に読むための形式的な接続口となる。ただし、本稿はこれらの計算効果がすべて同一構造であるとは主張せず、モナド、代数的効果、限制圏および時間付き意味論を置き換えない。向の理論が与えるのは、それらの内部意味論へ移る以前に、どこで対象世界が開かれ、どの情報が外延化によって失われるかを記録する境界幾何である。

\medskip
\noindent\textbf{キーワード：} 向、境界接続、投機、補償、残余、計算効果、非表現可能性、観測者、中観主義論理
\end{abstract}

\tableofcontents

\section{序論}

\subsection{向の存在意義をどこに置くか}

向の形式理論は、端を指定されない向、局所刻印、方向順序および刻印間の一致から事象と射を構成し、生経路圏から復元圏への商
\[
 Q_{\mathfrak D}:
 \mathcal P(\mathfrak D)\longrightarrow
 \mathcal C(\mathfrak D)
\]
を与えた\cite{ChibaDirection2026}。しかし自由圏と合同商による圏の復元だけでは、向を独立に導入する意義は十分ではない。圏論へ追加構造を付与すれば、時間、状態、確率、履歴および部分性の多くは個別に記述できるからである。

向の固有の役割は、内部の時間発展をすべて置き換えることではなく、次の二つを同時に可視化することにある。
\begin{enumerate}
 \item 過程表示が外延的な射へ送られる際に失われる細分、履歴および観測差。
 \item 選択した対象世界では終域を形成できない境界への半開接続。
\end{enumerate}

第一は商関手 $Q_{\mathfrak D}$ のファイバーに現れ、第二は非表現可能な境界方向前層に現れる。本稿は特に第二の構造を、将来の投機的・計算的解釈へ接続するための一般形式として検討する。

\subsection{内部時間と境界時間}

状態 $A_t,A_{t+1}$ が同じ対象世界の内部で形成されているなら、時間発展は
\[
 A_t\longrightarrow A_{t+1}
\]
という射で記述できる。必要なら射へ時間長、コスト、確率または履歴を付加すればよい。したがって向は、時間性一般のために必要なのではない。

一方、投機では、現在の内部世界からまだ検証されていない結果、継続または未来状態へ接続し、その後に成功、反証、取消しまたは修正として帰還する。この未来が既存の可能性空間の対象として既に形成されているなら、やはり射や関係で扱える。問題となるのは、未来側の候補全体が現在の観点における一つの内部対象へ回収されない場合である。この接面を
\[
 A\dline\dfix
\]
と記すことが、本稿における向と投機性の最小の接続である。

\subsection{主張の範囲}

本稿は次を主張する。
\begin{enumerate}
 \item 境界方向の正負プロファイルと部分的ペアリングは、内部世界と非対象化境界を結ぶ接続構造を与える。
 \item この接続構造により、恒等帰還、残余を伴う帰還、帰還不能を区別できる。
 \item 投機、補償および副作用は、一定の条件の下でこの境界構造の解釈例となり得る。
 \item 既存の計算効果意味論は不要にならず、向はそれらへ入る前の表示層として補完的に働く。
\end{enumerate}

本稿は「数学的対象はすべて閉じており、計算対象はすべて開いている」とは主張しない。また、例外、状態、遅延および投機が数学的に同型であるとも主張しない。開放性と外部性は、どの圏を内部世界として選ぶかに相対的である。

\section{向の形式理論から境界接続へ}

\subsection{復元圏と外延化}

向系 $\mathfrak D$ から復元される圏を $\mathcal C=\mathcal C(\mathfrak D)$ とする。生経路 $p,q\in\mathcal P(\mathfrak D)(A,B)$ が
\[
 Q_{\mathfrak D}(p)=Q_{\mathfrak D}(q)
\]
を満たすとき、両者は同じ外延的射を表すが、刻印の細分や向の乗換え表示は異なり得る。したがって $Q_{\mathfrak D}$ は、単に圏を作る写像ではなく、過程表示から外延的振舞いへの忘却と読める。

この忘却は時間量そのものを必ず消すわけではない。時間量を復元射のラベルとして保持する模型も可能である。重要なのは、何を合同で同一視し、何を射の不変量として残すかが、観測と意味論の選択に依存することである。

\subsection{境界方向前層}

理想刻印 $\omega$ に対し、内部対象 $A$ から境界へ向かう半開方向の集合を
\[
 \partial_\omega^+(A)
\]
とする。内部射 $f:A'\to A$ に沿う前合成によって
\[
 \partial_\omega^+:
 \mathcal C^{\mathrm{op}}\longrightarrow\Set
\]
を得る。$\partial_\omega^+$ が表現可能でないとき、すべての半開方向を一つの内部対象 $R$ への射
\[
 \Hom_{\mathcal C}(-,R)
\]
として回収できない。この非表現可能性が、現在の対象世界に相対的な境界性を与える。

反対向系または別途与えた帰還データから、境界側から内部対象 $B$ へ戻る負の方向の集合
\[
 \partial_\omega^-(B)
\]
を得る。本稿では、内部射 $g:B\to B'$ に沿って帰還先を移送できるという追加仮定の下で
\[
 \partial_\omega^-:
 \mathcal C\longrightarrow\Set
\]
を余前層として扱う。この共変性は向系の基本公理から自動的に従うとは限らず、境界接続を論じるための構造仮定である。

\subsection{拡大圏による表現}

境界方向の非表現可能性は、絶対的な記述不能性ではない。関手
\[
 j:\mathcal C\longrightarrow\mathcal D
\]
と境界対象 $\omega\in\mathcal D$ を持つ拡大圏では
\[
 \partial_\omega^+(A)
 \cong
 \Hom_{\mathcal D}(jA,\omega)
\]
と表現できる。このとき向は拡大圏内の射へ戻るが、$\omega$ が $j$ の本質像へ降下しないこと、および前層が $\mathcal C$ 内で表現可能でないことは残る。これは中観主義論理の境界方向意味論と同じ相対化である\cite{ChibaMadhyic2026}。

したがって向の役割は、圏論的表現を禁止することではない。内部化のためにどの意味論的拡大を行ったかを隠さず、拡大以前の観点では終域が所有されていなかったことを記録する点にある。

\section{境界接続子}

\subsection{正負方向の適合性}

外向き方向 $r^+\in\partial_\omega^+(A)$ と帰還方向 $r^-\in\partial_\omega^-(B)$ が常に一つの内部射を定めるとは限らない。両者が異なる境界成分、異なる完備化または両立しない観測条件を参照する可能性があるからである。

\begin{definition}[境界適合関係]
各 $A,B\in\mathcal C$ に対し
\[
 K_\omega(A,B)
 \subseteq
 \partial_\omega^+(A)\times\partial_\omega^-(B)
\]
を与える。$(r^+,r^-)\in K_\omega(A,B)$ のとき、正負の方向は $A$ から $B$ への帰還について\textbf{境界適合的}であるという。
\end{definition}

適合関係を明示することで、帰還不能を空の射や底対象へ直ちに同一視せずに保持できる。

\subsection{部分的境界ペアリング}

\begin{definition}[境界接続子]
境界接続子とは、各 $A,B\in\mathcal C$ に対する写像
\begin{equation}
 \kappa_{A,B}:
 K_\omega(A,B)
 \longrightarrow
 \Hom_{\mathcal C}(A,B)
 \label{eq:boundary-connector}
\end{equation}
の族である。部分写像の記法では
\[
 \kappa_{A,B}:
 \partial_\omega^+(A)\times\partial_\omega^-(B)
 \rightharpoonup
 \Hom_{\mathcal C}(A,B)
\]
と書く。
\end{definition}

$\kappa$ は、$\omega$ を内部対象とする射の合成ではない。境界を経由する正負の方向表示から、適合する場合に限って内部射を生成する追加構造である。

\begin{assumption}[接続子の自然性]
$f:A'\to A$、$g:B\to B'$ と $(r^+,r^-)\in K_\omega(A,B)$ に対し、移送された組
\[
 \bigl(\partial_\omega^+(f)(r^+),
       \partial_\omega^-(g)(r^-)\bigr)
\]
は $K_\omega(A',B')$ に属し
\begin{equation}
 \kappa_{A',B'}
 \bigl(\partial_\omega^+(f)(r^+),
       \partial_\omega^-(g)(r^-)\bigr)
 =g\circ\kappa_{A,B}(r^+,r^-)\circ f
 \label{eq:connector-naturality}
\end{equation}
を満たすと仮定する。
\end{assumption}

この自然性は、境界表示の選び方と内部の前後処理が整合することを要求する。具体的模型では、適合関係の移送安定性とともに検証しなければならない。

\subsection{ホロノミー}

\begin{definition}[境界ホロノミー]
$(r^+,r^-)\in K_\omega(A,A)$ に対し
\[
 \Hol_{A,\omega}(r^+,r^-)
 :=\kappa_{A,A}(r^+,r^-)
 \in\End_{\mathcal C}(A)
\]
と定める。
\end{definition}

これは向の形式理論で導入された自己境界ペアリングを、異なる帰還先を許す接続子の対角部分として捉え直したものである。$A\dline\dfix\dline A$ は、$\dfix$ を対象とする合成ではなく、この対角ペアリングの表示である。

\section{投機性への形式的フック}

\subsection{投機を何と区別するか}

未知であることだけでは、向を必要としない。未知量が既存の標本空間、探索空間または型の内部に置かれているなら、確率、関係、非決定性または部分写像で扱える。投機性に固有なのは、結果が未確定であることに加えて、現在の文脈が未来の結果を先取りして行動し、その後に検証または取消しを受けるという非対称性である。

本稿では投機の完全な計算体系を与えず、次の解釈だけを接続口として残す。
\begin{center}
\begin{tabular}{c p{0.60\textwidth}}
\toprule
境界データ & 投機的解釈 \\
\midrule
$r^+\in\partial_\omega^+(A)$
 & 現在の文脈から未確定な未来または継続へ向かう投機 \\
$r^-\in\partial_\omega^-(B)$
 & 検証、反証、結果の取込みまたは制御の復帰 \\
$(r^+,r^-)\in K_\omega(A,B)$
 & 投機と帰還が同じ境界事象について適合すること \\
$\kappa_{A,B}(r^+,r^-)$
 & 検証後に内部世界へ得られる外延的結果 \\
\bottomrule
\end{tabular}
\end{center}

この表は意味論の定義ではなく、将来の投機計算が境界接続子をどのように実現すべきかを示すインターフェイス仕様である。

\subsection{成功、取消しおよび未解決}

接続子により、少なくとも次を区別できる。
\begin{enumerate}
 \item $A$ から $B$ への成功：$(r^+,r^-)\in K_\omega(A,B)$ で、$\kappa_{A,B}(r^+,r^-)$ が結果射を与える。
 \item 完全な取消し：$B=A$ かつ $\Hol_{A,\omega}(r^+,r^-)=\id_A$ である。
 \item 痕跡を伴う帰還：$B=A$ だが $\Hol_{A,\omega}(r^+,r^-)\neq\id_A$ である。
 \item 未解決または帰還不能：$r^+$ に適合する $r^-$ が存在しない。
\end{enumerate}

これらは成功、失敗、例外、クラッシュと直ちに同一ではない。どの状態を成功と呼び、未解決を停止不能、例外または外部依存のどれとして読むかは、別途与える操作意味論に依存する。

\subsection{投機の時間性}

投機の時間性は、未来を一本の内部時間軸として先へ進むことだけにあるのではない。現在の文脈が、まだ内部対象として確定していない結果に依存して先行し、その依存が後から検証される点にある。この意味で
\[
 A\dline\dfix\dline B
\]
は、$A$ から $B$ までの時間経路そのものではなく、内部時間軸を成立させる前後に置かれた境界接続を示す。

将来の理論では、境界接続へ時刻、期限、予測根拠および検証証拠を付加できる。しかし、それらは向の一般理論の公理ではなく、投機的解釈の付加構造である。

\section{補償と残余}

\subsection{補償可能性}

\begin{definition}[補償可能な外向き方向]
$r^+\in\partial_\omega^+(A)$ が\textbf{補償可能}であるとは、ある $r^-\in\partial_\omega^-(A)$ が存在して
\[
 (r^+,r^-)\in K_\omega(A,A),
 \qquad
 \Hol_{A,\omega}(r^+,r^-)=\id_A
\]
となることをいう。この $r^-$ を $r^+$ の\textbf{完全帰還証拠}と呼ぶ。
\end{definition}

この定義は、境界を内部対象 $\omega$ とする射 $A\to\omega\to A$ を仮定しない。したがって向の非対象化条件と両立する。一方、$\omega$ を含む拡大圏または部分写像の圏が与えられた模型では、完全帰還証拠を部分逆、restriction retraction または補償操作として実現できる場合がある。限制圏は部分写像とその定義域をrestriction idempotentによって記述するため、この実現を検討する有力な比較対象である\cite{CockettLack2002}。ただし一般の境界接続子が自動的に限制圏へ還元されるとは主張しない。

\subsection{残余}

一般の圏では自己射から恒等射を減算できないため、残余を
\[
 \Hol_{A,\omega}(r^+,r^-)-\id_A
\]
という差として定義できない。

\begin{definition}[残余的帰還]
適合する組 $(r^+,r^-)\in K_\omega(A,A)$ が
\[
 \Hol_{A,\omega}(r^+,r^-)\neq\id_A
\]
を満たすとき、これを\textbf{残余的帰還}と呼び、自己射 $\Hol_{A,\omega}(r^+,r^-)$ 自体を帰還後の残余プロファイルと読む。
\end{definition}

数値的な残余量または「厚み」を定義するには、$\End(A)$ 上の距離、順序、ノルム、トレースまたは観測同値などの追加構造が必要である。本稿では残余を恒等性からの定性的な逸脱としてのみ扱う。

\subsection{三種類の開放性}

境界方向の有無と帰還の性質から、次の三状態を区別できる。
\begin{center}
\begin{tabular}{c p{0.62\textwidth}}
\toprule
状態 & 条件 \\
\midrule
境界閉鎖的
 & 関心のある正の境界方向が存在しない。 \\
補償可能に開放的
 & 正の境界方向は存在するが、各方向に恒等帰還を与える適合方向が存在する。 \\
残余的に開放的
 & ある方向が帰還不能であるか、すべての帰還が非恒等ホロノミーを誘導する。 \\
\bottomrule
\end{tabular}
\end{center}

これは「数学は厚み零、計算は厚み正」という絶対分類ではない。同じ対象も、純粋層、状態層、環境層のどれを内部に含めるかによって異なる状態へ分類される。

\section{副作用との接続}

\subsection{副作用も相対的である}

状態、入出力、例外および非決定性は、モナドやKleisli圏を用いて圏論的に内部化できる\cite{Moggi1991,PlotkinPower2002}。したがって副作用は絶対に射で書けない現象ではない。純粋計算の圏 $\mathcal C$ に対して効果を外部と見る場合には境界接続として現れ、効果を含むKleisli圏を内部世界として選べば通常の射として現れる。

向の理論が記録するのは、この内部化の前後である。すなわち、どの環境依存を境界へ押し出したか、どの拡大によって効果を内部射へ変えたか、および帰還後にどの差が残ったかを表示する。

\subsection{解釈例}

以下は同型対応ではなく、境界接続子の解釈候補である。
\begin{center}
\begin{tabular}{p{0.23\textwidth} p{0.44\textwidth} p{0.22\textwidth}}
\toprule
現象 & 境界接続としての読み & 限定 \\
\midrule
外部状態・入出力
 & 環境へ向かう外向き方向と、応答または更新後の帰還
 & 環境を基底圏の外部に置く場合 \\
投機的実行
 & 未検証結果への外向き方向と、検証または取消しによる帰還
 & 未来結果がまだ内部対象化されない場合 \\
トランザクション
 & 一時的な外部状態への方向と、commitまたはrollbackによる帰還
 & 原子的意味論を別途必要とする \\
例外
 & 通常継続から外れた境界と、handlerによる別の帰還先
 & 例外型を内部化すれば通常の射になる \\
遅延評価
 & 未評価計算への依存と、需要発生後の値の帰還
 & 逆時間方向とはみなさない \\
命令的代入
 & 状態環境を外部に置く場合の更新と帰還
 & 純粋な構文的代入とは区別する \\
\bottomrule
\end{tabular}
\end{center}

この表の意義は効果を新しく定義することではなく、異なる効果が「内部世界から一度外へ出て、結果または痕跡を伴って戻る」という表示を共有し得ることを示す点にある。

\subsection{投機と副作用は同一ではない}

投機は未来結果を先取りする認識的・時間的非対称性を持つ。副作用は環境または状態の変化であり、未来予測を伴わなくても生じる。両者が重なるのは、投機的実行が外部状態を変更し、その後の取消しで完全には元へ戻れない場合である。

したがって向の理論は投機と副作用を同一概念へ還元しない。両者が共有する境界接続形式と、各応用固有の意味論を分離する。

\section{観測者と残余の可視性}

\subsection{方向的lumpabilityとの区別}

方向的lumpabilityは、観測者による粗視化 $\pi_o$ の下で、同じ観測状態を表す刻印が同じ観測未来プロファイルを持つ条件である\cite{ChibaDirection2026,KemenySnell1976}。これは帰還が恒等であることとは別である。

完全に補償可能な過程でも、観測者が途中の情報を見れば異なる未来を区別できる場合がある。逆に、内部には非恒等な残余があっても、観測者がその成分を忘れるなら外からは恒等帰還に見える場合がある。

\subsection{観測的残余}

観測圏 $\mathcal E_o$ と関手
\[
 O_o:\mathcal C\longrightarrow\mathcal E_o
\]
を取る。

\begin{definition}[観測的に無残余な帰還]
適合する $(r^+,r^-)\in K_\omega(A,A)$ が観測者 $o$ に対して\textbf{無残余}であるとは
\[
 O_o\bigl(\Hol_{A,\omega}(r^+,r^-)\bigr)
 =\id_{O_o(A)}
\]
となることをいう。
\end{definition}

内部的には残余的であっても観測的には無残余となり得る。この相対性により、ログ、キャッシュ、プロファイル、外部送信など、どの痕跡を観測者が区別するかをモデルごとに指定できる。

\begin{proposition}
完全帰還
\[
 \Hol_{A,\omega}(r^+,r^-)=\id_A
\]
は任意の関手 $O_o$ に対して観測的に無残余である。
\end{proposition}

\begin{proof}
関手は恒等射を恒等射へ送るため
\[
 O_o(\Hol_{A,\omega}(r^+,r^-))
 =O_o(\id_A)
 =\id_{O_o(A)}
\]
である。
\end{proof}

逆は一般に成立しない。$O_o$ が忠実でない場合、異なる内部自己射が同じ観測射へ送られ得るからである。

\section{既存理論との関係}

\subsection{モナドと代数的効果}

モナド意味論は計算効果を型 $A$ から計算型 $TA$ への移行として内部化し、効果付き計算をKleisli射として扱う\cite{Moggi1991}。向の境界接続子は、この内部化を置き換えない。むしろ
\[
 \text{純粋層}
 \dline
 \text{効果境界}
 \dline
 \text{効果付き内部層}
\]
という意味論的拡大の表示を与える候補である。

具体的な比較には、境界方向からKleisli射への写像、接続子とKleisli合成の整合性、および非表現可能性がモナドによる自由代数化でどのように変化するかを示す必要がある。これらは未構成である。

\subsection{限制圏と部分写像}

限制圏は、部分写像 $f:A\to B$ に定義域を表すrestriction idempotent $\overline f:A\to A$ を対応させる\cite{CockettLack2002}。帰還不能や部分的な接続子をモデル化する際に有用であるが、理想刻印 $\omega$ が内部対象でない限り、境界方向 $A\dline\omega$ を限制圏の射 $A\to\omega$ と直接同一視できない。

比較を行うには、$\omega$ を含む拡大限制圏を構成するか、部分接続子 $\kappa$ の定義域 $K_\omega(A,B)$ をrestriction構造として実現する必要がある。本稿は可能性だけを指摘し、一般的な還元定理を主張しない。

\subsection{有向位相と高次構造}

有向代数的位相は非可逆な路と基本圏を扱い\cite{Grandis2003Directed}、合成的高次圏論は有向区間と高次整合性を型理論へ内部化する\cite{RiehlShulman2017}。これらの理論では点、型または有向路が既に与えられる。向の理論は、その前段階として刻印と一致から対象を形成し、さらに実現されない境界刻印を区別する。

ただし境界接続子の高次整合性は未定義である。異なる正負方向対が同じ内部射を生む理由、補償証拠間の変形、および残余の合成を扱うには、二次セルまたはホモトピー的構造が必要になる。

\section{中観主義論理および共観主義論理との接続}

\subsection{中観主義論理}

中観主義論理では、観点圏 $\mathcal C_i$、拡大圏 $\mathcal D$、解釈 $j_i$ および非降下境界 $\gamma$ から
\[
 \partial_\gamma^{(i)}(X)
 =\Hom_{\mathcal D}(j_iX,\gamma)
\]
を構成し、これが内部で表現可能でないことを外部方向条件とした\cite{ChibaMadhyic2026}。向の境界接続子は、この正方向プロファイルへ負方向と帰還ペアリングを追加する展望を与える。

正方向だけでは、外部へ向かう指示は記述できても、その指示が検証、反証または内部変化としてどのように戻るかは決まらない。接続子はこの帰還を追加構造として分離し、外部性と帰還効果を混同しない。

\subsection{共観主義論理}

共観主義論理は複数の観点から共通比較層へ向かう支持と非降下な比較余剰を扱う\cite{ChibaCotuition2026}。複数観点 $i$ が同じ境界 $\omega$ へ正方向を持つ場合、各観点の投機または探索を同じ境界上で比較できる可能性がある。

しかし、同じ境界表示を持つことから同じ帰還結果が従うわけではない。観点ごとの負方向、適合関係および接続子を比較する自然変換が必要である。この比較は、複数の投機主体、分散計算または異なる観測フレームを扱う将来課題となる。

\section{具体的研究課題}

\subsection{接続子の存在模型}

第一の課題は、$\partial_\omega^\pm$ と $\kappa$ が非自明に存在する具体例を与えることである。必要なのは、単に境界を対象として追加することではなく、次を同時に満たす模型である。
\begin{enumerate}
 \item 正方向前層が基底圏で非表現可能である。
 \item 負方向が共変移送を持つ。
 \item 適合関係が内部射の前後合成に安定である。
 \item 接続子が自然性\eqref{eq:connector-naturality}を満たす。
 \item 恒等帰還と残余的帰還の両方が存在する。
\end{enumerate}

自然数の無限遠方向は正方向の非表現可能性を与えるが、非自明な帰還接続子を自動的には与えない。計算的応用には、状態遷移系、トランザクション模型またはゲーム意味論的な境界を用いる必要があると予想される。

\subsection{接続子の圏論的位置づけ}

$\partial_\omega^+$ は前層、$\partial_\omega^-$ は余前層であり、$\kappa$ は両者の積からHom関手への部分的な評価に似ている。これをプロ関手、双加群、coendまたはChu型双対性のどれとして定式化すべきかは未決定である。

\begin{conjecture}[境界接続子のプロ関手表示]
適切な小ささ、自然性および適合関係の表現可能性を仮定すれば、境界接続子を二つのプロ関手の合成からHomプロ関手への射として表現できる具体的クラスが存在する。
\end{conjecture}

この予想が成立すれば、異なる境界の合成、接続子の随伴性および高次化を既存のプロ関手計算へ接続できる。

\subsection{投機計算の独立形式化}

投機を本格的に扱うには、少なくとも次のデータが必要である。
\begin{enumerate}
 \item 投機開始、予測根拠および依存する未確定値。
 \item 検証、反証、取消しおよび期限切れの区別。
 \item 投機中に行われた内部・外部効果の記録。
 \item 完全補償、部分補償および補償不能の合成則。
 \item 観測者ごとの残余の可視性。
\end{enumerate}

この体系は向の一般理論とは分離した独立論文として構成すべきである。向の側には、正負方向、適合関係、接続子およびホロノミーという型だけを残す。

\subsection{機械化}

境界接続子は部分構造を含むため、型理論または証明支援系での機械化が有用である。特に、適合関係の証拠を接続子の入力型へ含めれば
\[
 \kappa_{A,B}:
 \sum_{(r^+,r^-)}
 \operatorname{Compatible}(r^+,r^-)
 \longrightarrow
 \Hom_{\mathcal C}(A,B)
\]
と全域関数として記述できる。これにより、帰還不能を暗黙の未定義性ではなく、適合証拠の不在として扱える。

\section{結論}

向の理論の価値は、すべての時間過程を新しい原始概念へ置き換えることにはない。状態と終域が内部対象として形成された後の時間発展は通常の射で扱える。向が必要になるのは、過程が外延的な射へ商される以前の表示を保持するとき、および選択した対象世界では終域を形成できない境界へ接続するときである。

本稿は、この境界接続を
\[
 \boxed{
 \partial_\omega^+(A)
 \times
 \partial_\omega^-(B)
 \supseteq K_\omega(A,B)
 \xrightarrow{\ \kappa_{A,B}\ }
 \Hom_{\mathcal C}(A,B)
 }
\]
として整理した。対角成分は境界ホロノミーを与え、恒等帰還、残余的帰還および帰還不能を区別する。これにより
\[
 A\dline\dfix\dline B
\]
は、$\dfix$ を通過対象とする射の合成ではなく、非対象化境界を介して内部射が生成される接続図式として読まれる。

投機的解釈では、正方向は未確定な未来への先行、負方向は検証または復帰、接続子は検証後の内部結果を表す。副作用との接続では、非恒等ホロノミーが帰還後の痕跡を表す。ただしこれらは境界構造の解釈であり、計算効果の既存意味論を置き換えない。モナド、限制圏、時間付き意味論および有向高次構造との厳密な比較は、具体模型を構成した後の課題である。

したがって向の理論の展望は、圏から何を構成できるかだけでなく、圏へ外延化する前に何が存在し、その外延化で何が失われ、境界からの帰還がどの残余を残すかを記述することにある。この接続構造が投機性の時間意味論へ至る形式的なフックとなる。

\bibliographystyle{plain}
\bibliography{ref}

\end{document}